% ============================================================================
%  GENERATED by scripts/gen-paper-tex.mjs from apps/web/lib/papers.ts
%  DO NOT EDIT BY HAND. Edit the manifest and regenerate:
%
%      node scripts/gen-paper-tex.mjs
%
%  The HTML at /papers/hardwood-selection-and-cost-framework-gta and this document are the same manuscript.
%  That is the point: a hand-written .tex beside a manifest-driven page is two
%  documents that agree until the day they do not.
%
%  To publish the PDF:
%      1. compile this file (pdflatex, twice, for the table of contents)
%      2. mkdir -p apps/web/public/papers && mv <the-new>.pdf apps/web/public/papers/
%      3. git add apps/web/public/papers   -- in the same commit
%  The download button appears on its own; pdfIsPublished() checks the file.
% ============================================================================
\documentclass[11pt,a4paper]{article}
\usepackage[utf8]{inputenc}
\usepackage[T1]{fontenc}
\usepackage[margin=2.4cm]{geometry}
\usepackage{booktabs}
\usepackage{longtable}
\usepackage{array}
\usepackage{ragged2e}
\usepackage{enumitem}
\usepackage{xcolor}
\usepackage{mdframed}
\usepackage{titlesec}
\usepackage{fancyhdr}
\usepackage[hidelinks]{hyperref}

\definecolor{ecodark}{RGB}{16,20,28}
\definecolor{ecowood}{RGB}{145,95,45}
\definecolor{ecoaccent}{RGB}{0,100,120}
\definecolor{ecolite}{RGB}{250,247,242}

\titleformat{\section}{\Large\bfseries\color{ecodark}}{\thesection}{0.7em}{}
\titleformat{\subsection}{\large\bfseries\color{ecoaccent}}{}{0em}{}
\setlist[itemize]{leftmargin=1.4em,itemsep=0.25em,topsep=0.4em}
\setlist[enumerate]{leftmargin=1.6em,itemsep=0.25em,topsep=0.4em}

\newmdenv[
  linecolor=ecowood,linewidth=2pt,topline=false,bottomline=false,rightline=false,
  backgroundcolor=ecolite,leftmargin=0pt,rightmargin=0pt,
  innerleftmargin=1em,innertopmargin=0.7em,innerbottommargin=0.7em,skipabove=1em,skipbelow=1em
]{callout}

\pagestyle{fancy}
\fancyhf{}
\fancyhead[L]{\footnotesize\color{ecodark}EcoWoods Hardwood Flooring \textbullet{} Toronto}
\fancyhead[R]{\footnotesize\color{ecodark}The Intelligent Homeowner's Decision Framework v1.0}
\fancyfoot[C]{\footnotesize\thepage}
\renewcommand{\headrulewidth}{0.4pt}

\begin{document}
\begin{titlepage}
\vspace*{3cm}
\begin{flushleft}
{\Huge\bfseries\color{ecodark}The Intelligent Homeowner's Decision Framework}\\[0.6em]
{\LARGE\color{ecoaccent}How to choose hardwood that performs, appreciates, and never becomes a liability}\\[2em]
{\large A decision framework for hardwood in the Greater Toronto Area: installed cost ranges, species hierarchy by Janka hardness, the solid-versus-engineered decision tree, and a checklist for evaluating any installer.}\\[2.5em]
\rule{\textwidth}{0.6pt}\\[0.8em]
\textbf{Version} 1.0 \quad\textbullet\quad \textbf{Published} 2026-08-01 \quad\textbullet\quad \textbf{Reading time} 8 minutes\\[0.4em]
\textbf{Audience} Homeowners planning a renovation, designers, realtors\\[0.4em]
\textbf{Topics} Cost · Species selection · Janka hardness · Resale · Installer evaluation\\[1.2em]
\small\textit{The current edition of this paper is published as HTML at }\texttt{ecowoods.ca/papers/hardwood-selection-and-cost-framework-gta}\textit{. Where the two disagree, the page is correct.}
\end{flushleft}
\end{titlepage}

\tableofcontents
\newpage

\section*{Summary}
\addcontentsline{toc}{section}{Summary}
Hardwood is one of the few renovations that reliably returns more than it costs --- when the physics and the process are both handled correctly. This paper covers what installation actually costs in the GTA, how to choose a species, how the substrate decides the product, and the six questions that separate a competent installer from a cheap one.

\section{Hardwood as a renovation investment}
A correctly specified and installed hardwood floor is one of the few renovations that reliably returns more than it costs. It is the strongest visual signal to buyers, and a fixed-price installation removes the change-order risk that makes buyers nervous about a renovated home.

\begin{itemize}
  \item 70--100\%+ cost recovery on resale in the GTA
  \item Strongest visual and emotional signal to buyers
  \item Removes the primary buyer fear: future change orders and failure
\end{itemize}

\section{Installed cost in the GTA}
These ranges assume professional installation, correct product, and the full protocol. Bargain pricing almost always means skipped moisture testing, shorter acclimation, or subcontracted labour.


\begin{center}
\footnotesize
\begin{longtable}{p{8.00cm}p{8.00cm}}
\caption*{\small\itshape Fully installed pricing, Greater Toronto Area}\\
\toprule
\textbf{Scope} & \textbf{Range} \\
\midrule
\endhead
Fully installed, average (GTA market) & ≈ \$13 / sq ft \\
Fully installed, typical range (GTA market) & \$8 -- \$18 / sq ft \\
Screen and recoat (Screen \& Recoat) & \$2.50--\$4.00 / sq ft \\
Full sand and finish (Full Sand \& Finish) & \$4.75--\$7.50 / sq ft \\
Premium new install (wide-plank, oil finishes, stairs) & \$11.00--\$18.00 / sq ft \\
\bottomrule
\end{longtable}
\end{center}
\section{Species hierarchy}
Wide-plank European white oak currently dominates both aesthetics and resale signalling across the Greater Toronto Area. Hardness is only one variable --- stability, grain character and finish performance matter equally.


\begin{center}
\footnotesize
\begin{longtable}{p{5.33cm}p{5.33cm}p{5.33cm}}
\caption*{\small\itshape Janka hardness and role in the GTA}\\
\toprule
\textbf{Species} & \textbf{Janka} & \textbf{Role} \\
\midrule
\endhead
White oak / European oak & ≈1360 & Current aesthetic and resale sovereign \\
Hard maple & 1450 & High-traffic workhorse \\
Red oak (northern) & ≈1290 & Traditional default \\
Hickory & 1820 & Extreme durability \\
Black walnut & 1010 & Luxury accent \\
\bottomrule
\end{longtable}
\end{center}
\section{Solid versus engineered --- the decision tree}
The substrate decides the product, not the budget and not the preference.

\begin{enumerate}
  \item Is the substrate plywood over joists? → Solid is possible.
  \item Is the substrate concrete, condo slab, or radiant? → Engineered is required.
  \item Is the home subject to large seasonal RH swings? → Engineered preferred.
  \item Does the client want maximum future refinishing cycles? → Solid, only if the substrate allows.
\end{enumerate}

\begin{callout}
\textbf{\color{ecowood}Rule}\\[0.3em]
We specify. We never sell what the house cannot support for decades. The correct product is the one that will still look and perform correctly in twenty years.
\end{callout}

\section{What a fixed price actually protects}
A fixed price is not a marketing slogan. It is proof that the company has already performed the due diligence most installers skip --- because it forces them to own every decision before work begins.

\begin{itemize}
  \item The number written on the estimate is the number that will be paid
  \item No change orders for "unforeseen conditions" that should have been tested on day one
  \item Removes the renovation anxiety that destroys buyer confidence
\end{itemize}

\section{How to evaluate any installer}
Any "no" answer is a red flag.

\begin{enumerate}
  \item Do they moisture-test at the free estimate, and document the readings?
  \item Do they require a minimum 72-hour acclimation in the actual space?
  \item Is the price fixed in writing, with no open-ended change-order language?
  \item Are the installers salaried employees, or day-labour subcontractors?
  \item Will they refuse the job if the substrate or conditions are wrong?
  \item Is there a true lifetime workmanship warranty in the contract?
\end{enumerate}

\section{Homeowner action plan}
\begin{enumerate}
  \item Request a free in-home estimate that includes moisture testing and samples.
  \item Insist on a written fixed-price proposal with the full protocol listed.
  \item Confirm that only salaried artisans will perform the work.
  \item Verify the lifetime workmanship warranty language before signing.
  \item Keep indoor RH between 35\% and 55\% year-round after installation.
\end{enumerate}

\end{document}
